\input{../tex-inputs/latex-leseansicht-vorspann}

\section[Gustav Schwarzkopf an Arthur Schnitzler, 13. 8. 1902]{L04307 Gustav Schwarzkopf an Arthur Schnitzler, 13. 8. 1902}
\nopagebreak\mylabel{L04307v}
\rehead{ }\normalsize\beginnumbering\briefempfaengerindex{Schnitzler, Arthur@\textsc{Schnitzler, Arthur}!zzzSchwarzkopf, Gustav@\emph{von Gustav Schwarzkopf}!1902-08-131@{13. 8. 1902}|(be}
\toendnotes[C]{\smallbreak\pagebreak[2]}
\correspDesc{Versand  durch Gustav Schwarzkopf am 13. 8. 1902 in Strobl
\newline{}Erhalt  durch Arthur Schnitzler im Zeitraum [14. 8. 1902 – 18. 8. 1902?] in Wien}\toendnotes[C]{\smallbreak}
\Standort{CUL, Schnitzler, B 96a.}
\physDesc{Brief, 1 Blatt, 2 Seiten, 736 Zeichen
\newline{}Handschrift: schwarze Tinte, deutsche Kurrent}\toendnotes[C]{\smallbreak}
\pstart
           \raggedleft{}{\pb}\textsc{Strobl\oindex{Strobl@\textbf{Strobl}, \emph{Verwaltungsgebiet}|pw}}{ }13. 8. 02\pend
           \vspace{0.5em}
\pstart
           Lieber Arthur, ich gratulire, Ihnen herzlichſt und bitte Sie auch
               der Mama\pwindex{Schnitzler, Olga 17.\,1.\,1882 Wien – 13.\,1.\,1970 Lugano@\textsc{Schnitzler, Olga} (17.\,1.\,1882 Wien – 13.\,1.\,1970 Lugano), \emph{Schauspielerin, Sängerin}|pwv} des jungen Mannes\pwindex{Schnitzler, Olga 17.\,1.\,1882 Wien – 13.\,1.\,1970 Lugano@\textsc{Schnitzler, Olga} (17.\,1.\,1882 Wien – 13.\,1.\,1970 Lugano), \emph{Schauspielerin, Sängerin}|pwv} meine herzlichſten
               Glückwünſche zu melden. Ich hoffe ihn bald zu{ }ſehen, den Jüngling\pwindex{Schnitzler, Heinrich 9.\,8.\,1902 Hinterbrühl – 12.\,7.\,1982 Wien@\textsc{Schnitzler, Heinrich} (9.\,8.\,1902 Hinterbrühl – 12.\,7.\,1982 Wien), \emph{Regisseur, Schauspieler}|pwv}, – hat er{ }ſchon einen Namen? – Da
               ich die Abſicht habe, morgen nach Wien\oindex{Wien@\textbf{Wien}, \emph{Verwaltungsgebiet}|pw} zu fahren. Eigentlich wollte ich{ }ſchon Montag fort, hab
               mich aber – Schwachheit, Dein Name iſt Ma{\geminationn} - noch
               zurückhalten laſſen. Mein Katarrh iſt nicht beſſer geworden, er hat{ }ſich {\pb}nur verschli{\geminationm}ert, was leicht zu erklären iſt, da ich unter{ }ſieben Tagen fünf ganz verregnete
               gehabt habe, die überdies noch eine faſt winterliche Temperatur aufweiſen.\pend
           
\pstart
           Mit meinen Ausflügen habe ich in dieſer Saiſon kein Glück.\pend
           
\pstart
           Auch baldiges Wiederſehen!{\\[\baselineskip]}Herzlichſt{\\[\baselineskip]}Ihr{\\[\baselineskip]}\spacefill\mbox{Gustav.}\pend
           \leftskip=0em{}
\pstart
           \noindent{}Grüße an Frl. \textsc{Liesel\pwindex{Steinrück, Elisabeth 19.\,11.\,1885 – 7.\,4.\,1920 Partenkirchen@\textsc{Steinrück, Elisabeth} (19.\,11.\,1885 – 7.\,4.\,1920 Partenkirchen)|pw}} und Herrn \textsc{Paul\pwindex{Marx, Paul 21.\,7.\,1879 Wien – 30.\,10.\,1956 ebd.@\textsc{Marx, Paul} (21.\,7.\,1879 Wien – 30.\,10.\,1956 ebd.), \emph{Regisseur, Schauspieler}|pw}}.\pend
           \selectlanguage{ngerman}\endnumbering\briefempfaengerindex{Schnitzler, Arthur@\textsc{Schnitzler, Arthur}!zzzSchwarzkopf, Gustav@\emph{von Gustav Schwarzkopf}!1902-08-131@{13. 8. 1902}|)be}\mylabel{L04307h}
\begin{anhang}
\end{anhang}\newcommand{\dateiname}{L04307}\newcommand{\titel}{Gustav Schwarzkopf an Arthur Schnitzler, 13. 8. 1902}\newcommand{\editorInnen}{Herausgegeben von Jahnke, SelmaMüller, Martin Anton}\input{../tex-inputs/latex-leseansicht-abspann}
